\subsection{Hypothetical sales-pipeline mathematics for a paid studio}

To evaluate the commercial mechanics of a technical studio, consider a hypothetical experimental admissions cohort for Bayes Institute:
\begin{itemize}
\item \textbf{Inquiry Volume:} 40 eligible software engineers request detailed curriculum information;
\item \textbf{Diagnostic Completion:} 24 applicants complete the mandatory 45-minute technical benchmark assessment (60.0\% completion rate);
\item \textbf{Admissions Consultation:} 16 qualified applicants attend a 20-minute structured fit conversation (66.7\% conversion);
\item \textbf{Offer Extension:} 10 applicants satisfy all technical and scheduling criteria and receive an offer of admission (62.5\% qualification rate);
\item \textbf{Paid Enrollment:} 7 applicants pay tuition and enroll (70.0\% offer-to-enrollment conversion).
\end{itemize}
Overall funnel conversion from initial inquiry to paid enrollment is:
\begin{equation}
\text{Funnel Conversion} = \frac{7}{40} = 17.5\%.
\end{equation}

At a tuition fee of ₹30,000 per student, gross tuition receipts total:
\begin{equation}
7 \times ₹30,000 = ₹2,10,000 \quad \text{(₹2.1 lakh)}.
\end{equation}
The fully loaded operational economics of this cohort are modeled as follows:
\begin{itemize}
\item \textbf{Direct Delivery Costs:} Dedicated cloud laboratory instances, grading honoraria for junior code reviewers, and software tooling = ₹85,000;
\item \textbf{Cash Acquisition Spending:} Targeted developer community sponsorships and computational benchmarking distribution = ₹15,000;
\item \textbf{Founder Commercial Effort:} 25 hours spent on marketing, diagnostic grading, and admissions interviews. At an internal shadow planning wage of ₹1,000 per hour, this represents ₹25,000 of founder commercial capital.
\end{itemize}
Net contribution after direct delivery and acquisition costs is calculated as:
\begin{equation}
\begin{aligned}
\Pi_{\text{pre-delivery}} &= ₹2,10,000 - ₹85,000 \\
&\quad - ₹15,000 - ₹25,000 \\
&= ₹85,000.
\end{aligned}
\end{equation}

However, the analytical accounting discipline developed in Chapter 17 reveals a critical vulnerability. If delivering the 8-week studio requires 100 additional hours of unpriced founder instruction, capstone grading, and student mentoring, that delivery labor carries an economic opportunity cost of:
\begin{equation}
100 \text{ hours} \times ₹1,000/\text{hour} = ₹1,00,000 \quad \text{(₹1.0 lakh)}.
\end{equation}
Accounting for this instructional labor turns the cohort's true economic contribution negative:
\begin{equation}
\Pi_{\text{true}} = ₹85,000 - ₹1,00,000 = -₹15,000.
\end{equation}
A high conversion rate (70\% from offer to close) cannot compensate for an unviable delivery model. Scaling requires standardizing rubrics and delegating code grading to trained teaching assistants.

\subsection{A college proposal can look attractive and still be unsustainable}

Consider an enterprise partnership proposal where a private engineering university offers Bayes Institute ₹4,00,000 (₹4 lakh) to deliver an advanced AI elective for 120 undergraduate students across a 14-week semester.

The preliminary project budget appears highly profitable:
\begin{itemize}
\item \textbf{Gross Contract Value:} ₹4,00,000;
\item \textbf{Direct Instructional Costs:} Mentors, lab assistants, and student materials = ₹2,10,000 (₹2.1 lakh);
\item \textbf{Commercial and Proposal Labor:} Faculty coordination and curriculum mapping = ₹40,000;
\item \textbf{Allocated Administrative Overhead:} Software licenses and proctoring = ₹30,000;
\item \textbf{Nominal Accounting Profit:} ₹4,00,000 - ₹2,10,000 - ₹40,000 - ₹30,000 = ₹1,20,000 (₹1.2 lakh).
\end{itemize}

However, rigorous contract examination reveals severe operational and cash flow hazards:
\begin{enumerate}
\item \textbf{Cash Disbursement Timing:} University disbursement terms dictate payment 120 days following final semester grade submission. Bayes Institute must finance ₹2.8 lakh of instructional costs out of cash reserves for nearly 6 months, incurring ₹15,000 in financing costs;
\item \textbf{Scheduling Conflicts:} The university timetable mandates 6 live lecture hours weekly on Tuesday and Thursday mornings: directly conflicting with the founder's full-time employment;
\item \textbf{Uncapped Grading Demands:} The contract lacks proctoring bounds, requiring faculty to grade 120 individual capstone projects manually without university teaching assistant support.
\end{enumerate}
Despite an apparent 30\% profit margin, executing this contract under standard delivery terms would jeopardize the founder's employment and deplete operating cash. The contract must be restructured into a weekend masterclass format with university-funded teaching assistants, or declined entirely.

\subsection{An enterprise opportunity should be evaluated at realistic probability}

An enthusiastic email from a corporate engineering manager stating that ``your evaluation diagnostic looks incredible, we definitely want to roll this out across our 200 developers'' is not a ₹10 lakh contract. It is an unverified expression of interest.

Bayes Institute evaluates enterprise deals through a calibrated milestone framework:
\begin{itemize}
\item \textbf{Gate 1: Problem Confirmation (10\% Probability):} Urgent evaluation failure documented and acknowledged by frontline engineering leads;
\item \textbf{Gate 2: Economic Buyer Access (25\% Probability):} Introduction established with the VP of Engineering or CTO controlling departmental budget;
\item \textbf{Gate 3: Scope and Price Authorization (50\% Probability):} Written statement of work and commercial pricing approved by the economic buyer;
\item \textbf{Gate 4: InfoSec and Procurement Clearance (75\% Probability):} Vendor security questionnaire approved and Master Services Agreement redlines completed;
\item \textbf{Gate 5: Purchase Order Executed (90\% Probability):} Formal corporate purchase order issued with binding payment dates;
\item \textbf{Gate 6: Cash Collected (100\% Probability):} Electronic funds transfer cleared into Bayes Institute's operating bank account.
\end{itemize}
Weighted pipeline forecasting uses these calibrated probabilities to model projected cash collections, preventing premature capital expenditure against informal verbal promises.

\subsection{Bayes's 1st sales system can fit in a spreadsheet}

Before investing in expensive enterprise customer relationship management software, Bayes Institute's initial sales engine operates with complete rigor inside a single, version-controlled spreadsheet. The system maintains 3 dedicated tabs for D2C Learners, College Partnerships, and Enterprise Services.

Every prospect record captures 12 mandatory operational fields:
\begin{enumerate}
\item Account or Applicant Name;
\item Source Channel and Decision Trigger;
\item Verified Problem Episode and Current Workaround;
\item Economic Buyer and Evaluator Contacts;
\item Diagnostic Assessment Score and Skill Deficits;
\item Current Pipeline Gate and Supporting Artifact;
\item Next Buyer-Owned Action and Target Date;
\item Documented Objections and Required Evidence;
\item Quoted Contract Value and Baseline Price;
\item Payment Schedule and Bank Disbursement Dates;
\item Committed Founder Instructional Hours;
\item Win/Loss Classification and Root Cause Analysis.
\end{enumerate}
Attaching the final signed SOW directly to the customer record ensures that operational delivery aligns with commercial representations.

\subsection{The founder's 1st 30-day sales experiment}

To test commercial demand without jeopardizing full-time employment, the founder executes a tightly bounded 30-day sales experiment:
\begin{itemize}
\item \textbf{Week 1 (Offer Design and Hypothesis):} Define a single testable offer: an 8-week Applied Retrieval Evaluation Studio priced at ₹30,000 for 10 working engineers. Author explicit prerequisite criteria, workload bounds (12 hrs/week), and syllabus rubrics. Conduct 5 discovery interviews with software engineers who recently deployed production RAG applications;
\item \textbf{Week 2 (Proof Artifact and Outbound):} Publish a technically substantive benchmarking report detailing retrieval evaluation failure modes. Distribute the report across developer channels alongside an invitation to complete a 45-minute diagnostic assessment;
\item \textbf{Week 3 (Discovery and Qualification):} Administer diagnostics to 20 applicants. Conduct 20-minute fit conversations with qualified candidates, documenting alternative tools and primary objections. Extend formal enrollment offers to 10 qualified applicants;
\item \textbf{Week 4 (Commitment and Delivery Audit):} Observe actual enrollment payments. Record fully loaded acquisition hours and cash expenditures. Audit teaching capacity, verifying that code review turnaround workflows are operational before cohort kickoff.
\end{itemize}
If zero applicants pay tuition, the experiment has succeeded in falsifying the initial commercial hypothesis at minimal cost. The founder revises the problem statement or target persona rather than wasting months building an unwanted curriculum.
